\subsection{Generated, Cyclic, and Finitely Generated Modules}

\begin{definition} \label{def:generated-submodule} \hlt{Generated Submodule}
\begin{enumerate}[label=(\roman*)]
\setlength{\itemsep}{0pt}
\item For $S\subseteq M$, \hlt{submodule generated by $S$} is $\langle\!\langle S\rangle\!\rangle_R \defeq \left\{\sum\limits_{j=1}^{n} r_js_j:n\geq1,\ r_j\in R,\ s_j\in S\right\}$; $\langle\!\langle\emptyset\rangle\!\rangle_R=\{0\}$.
\item A subset $S\subseteq M$ \hlt{spans} or \hlt{generates} $M$ if $\langle\!\langle S\rangle\!\rangle_R=M$.
\end{enumerate}
\end{definition}

\begin{proposition} \label{prop:generated-submodule-smallest} \hlt{Generated Submodule Is Smallest}\\
For every $S\subseteq M$, $\langle\!\langle S\rangle\!\rangle_R = \bigcap\limits_{\substack{N\leq M\\ S\subseteq N}}N$.
\begin{proof}
Let $G\defeq\langle\!\langle S\rangle\!\rangle_R$. By construction, $S\subseteq G$ and the submodule test gives $G\leq M$.\\
If $N\leq M$ and $S\subseteq N$, then $\sum\limits_{j=1}^{n}r_js_j\in N$, where $r_j\in R,\ s_j\in S$, so $G\subseteq N$.\\
\hlt{Structural Commentary}: Generation means closing a set under the operations the module axioms require.\\
\hlt{Key Insight}: The generated submodule adds no elements except those forced by finite $R$-linear combination.
\end{proof}
\end{proposition}

\begin{definition} \label{def:cyclic-and-finitely-generated-modules} \hlt{Cyclic and Finitely Generated Modules}
\begin{enumerate}[label=(\roman*)]
\setlength{\itemsep}{0pt}
\item For $v\in M$, the \hlt{cyclic submodule generated by $v$} is $\langle\!\langle v\rangle\!\rangle_R=Rv=\{rv:r\in R\}$.
\item $M$ is \hlt{finitely generated} if $M=\langle\!\langle v_1,\ldots,v_n\rangle\!\rangle_R$ for some $v_1,\ldots,v_n\in M$.
\item $M$ is \hlt{cyclic} if $M=Rv$ for some $v\in M$.
\end{enumerate}
\end{definition}

\begin{example} \label{ex:finitely-generated-module-with-non-fg-submodule} \hlt{A Finitely Generated Module with a Non-Finitely Generated Submodule}\\
Let $R=\mathbb{F}[x_1,x_2,\ldots]$. Then $R=\langle\!\langle1\rangle\!\rangle_R$ as an $R$-module.\\
Set $I\defeq (x_1,x_2,\ldots)=\{f\in R:f(0,0,\ldots)=0\}$. Then $I\trianglelefteq R$, hence $I\leq R$, but $I$ is not finitely generated.\\
If $I=\langle\!\langle p_1,\ldots,p_n\rangle\!\rangle_R$, choose $k$ such that $x_k$ occurs in no $p_i$. Let $\varphi_k:R\to\mathbb{F}[x_k]$ send $x_k\mapsto x_k$ and $x_j\mapsto0$ for $j\neq k$. Since $\varphi_k(p_i)=0$ for all $i$, every element of $\langle\!\langle p_1,\ldots,p_n\rangle\!\rangle_R$ maps to $0$, but $\varphi_k(x_k)=x_k\neq0$, contradiction.
\end{example}

\begin{definition} \label{def:module-linear-independence} \hlt{Linear Independence in a Module}
\begin{enumerate}[label=(\roman*)]
\setlength{\itemsep}{0pt}
\item Subset $S\subseteq M$ is \hlt{linearly independent over $R$} if,\\
$\forall$ distinct $v_1,\ldots,v_n\in S$, $r_1v_1+\cdots+r_nv_n=0$ $\Longrightarrow$ $r_1=\cdots=r_n=0$.
\item A subset $S\subseteq M$ is \hlt{linearly dependent} if it is not linearly independent.
\end{enumerate}
\end{definition}

\begin{example} \label{ex:module-dependence-need-not-solve-for-vector} \hlt{Dependence Without Solving for One Vector}\\
In the $\mathbb{Z}$-module $\mathbb{Z}$, $3\cdot2-2\cdot3=0$. Thus $\{2,3\}$ is linearly dependent.\\
However, $2\notin 3\mathbb{Z}$; $3\notin 2\mathbb{Z}$, so dependence need not express one element as a scalar multiple of another.
\end{example}

\begin{proposition} \label{prop:integral-domain-scaling-preserves-independence} \hlt{Scaling Independent Sets over a Domain}\\
Assume $R$ is an integral domain. If $\{v_1,\ldots,v_n\}$ is $R$-linearly independent and $0\neq r\in R$, then $\{rv_1,\ldots,rv_n\}$ is $R$-linearly independent.
\begin{proof}
Let $a_1,\ldots,a_n\in R$ and suppose $a_1(rv_1)+\cdots+a_n(rv_n)=0$. Then $(ra_1)v_1+\cdots+(ra_n)v_n=0$.\\
Independence gives $ra_i=0$ for all $i$. Since $R$ is a domain and $r\neq0$, $a_i=0$ for all $i$.\\
\hlt{Structural Commentary}: Domains preserve independence under nonzero scaling.\\
\hlt{Key Insight}: The failure of this statement is exactly zero-divisor behavior.
\end{proof}
\end{proposition}
